\textbf{Source metadata: author (MetaInlines)}

Owner-directed SplitZero research continuation

\textbf{Source metadata: date (MetaInlines)}

13 September 2026

\textbf{Source metadata: lang (MetaInlines)}

en

\textbf{Source metadata: title (MetaInlines)}

SGA trace, the retained diagonal, and the arithmetic period determinant

\subsection{1. The completed calculation and its role}\label{endpoint-sga_trace_period--the-completed-calculation-and-its-role}

This continuation uses the \textbf{actual inline text in the supplied SGA 5 master}, especially Exposé III, §6.8, together with the preceding polynomial-exponential construction. It calculates, rather than postulates, the relative duality and diagonal class of the current cyclic arithmetic module. That class then evaluates the determinant of the entire exponential period matrix.

The principal formulas are

ENDPOINTSOURCE62EECF4C17912CE87EEDAFB1SLOT0END

where every symbol and map is specified below. In particular,

ENDPOINTSOURCE62EECF4C17912CE87EEDAFB1SLOT1END

This evaluates a previously unspecified analytic determinant. Its coefficient derivatives recover \textbf{all the first q power traces of the original companion operator}, and hence its whole characteristic polynomial, including multiplicities. The original theta metric is not changed. Its comparison with this calculated determinant is carried through in §§10--12.

These are explicit applications and extensions of classical residue/trace and irregular-period methods; no global priority claim is made. SGA 5 supplies the residue/trace operation, not a Hermitian positivity theorem for the theta quotient. Bloch--Esnault study determinants and periods for irregular connections, including the exp(f) case; their abstract was inspected, but their full article was not successfully retrieved in this execution. The proof here is supplied in full at the types used here.

\subsection{2. Source intake and mathematical categories}\label{endpoint-sga_trace_period--source-intake-and-mathematical-categories}

The nine uploaded files are LaTeX masters. SGA 5 contains 15,818 lines of substantive inline material. The other masters have differing amounts of inline material and refer to further component files; for example the SGA 4½ master consists of its preamble and eight component inputs. The absent component contents are not counted as read. \texttt{SOURCE\_REVIEW.md} and \texttt{checks/source-inventory.json} give the exact reading and dependency inventory.

The key supplied statement is SGA 5, Exposé III, (6.8.5): the residue of a function multiplied by the regular-immersion class equals the trace of its multiplication endomorphism on a finite flat algebra. The preceding paragraph constructs the class from the conormal differential. The source explicitly records a qualification about the proof/reference of this compatibility. We retain that qualification and prove its one-variable instance directly below.

The appendix works with coherent complexes over a noetherian scheme, with the indicated finite Tor-dimension, properness, and Tor-independence hypotheses. Here the finite map is finite free, so properness and flatness are actual properties of the displayed algebra. Its dual is computed directly as a finite module dual. No general six-operation formalism on semiring schemes is assumed.

Keep the programme's one scalar construction

ENDPOINTSOURCE62EECF4C17912CE87EEDAFB1SLOT2END

Its support observation is ENDPOINTSOURCE62EECF4C17912CE87EEDAFB1SLOT3END and ENDPOINTSOURCE62EECF4C17912CE87EEDAFB1SLOT4END. The pair ENDPOINTSOURCE62EECF4C17912CE87EEDAFB1SLOT5END retains the whole scalar. For every ring map ENDPOINTSOURCE62EECF4C17912CE87EEDAFB1SLOT6END, its split lift fixes ENDPOINTSOURCE62EECF4C17912CE87EEDAFB1SLOT7END and sends ENDPOINTSOURCE62EECF4C17912CE87EEDAFB1SLOT8END to ENDPOINTSOURCE62EECF4C17912CE87EEDAFB1SLOT9END. In particular

ENDPOINTSOURCE62EECF4C17912CE87EEDAFB1SLOT10END

remains the original arithmetic square with an infinite left target. All calculations below first specify coefficient maps and then lift them fibrewise through the original support reconstruction. The ordinary coherent realization is connected to that richer object by p, not substituted for the complete split object.

At a label ℓ, an R-linear map f lifts as (ℓ,x) ↦ (ℓ,f(x)); a coefficient zero remains (ℓ,0). External absence maps to absence. Kernels and quotients below are the original coefficient kernels and quotients before that reconstruction. No new scalar zero is adjoined for a diagonal homology degree or a deformation parameter.

\subsection{3. The arithmetic object that is retained}\label{endpoint-sga_trace_period--the-arithmetic-object-that-is-retained}

Fix the original finite packet h, with complete orders, and tensor degree k. Retain

ENDPOINTSOURCE62EECF4C17912CE87EEDAFB1SLOT11END

The source and actual weighted inclusion are

ENDPOINTSOURCE62EECF4C17912CE87EEDAFB1SLOT12END

The retained source identities are

ENDPOINTSOURCE62EECF4C17912CE87EEDAFB1SLOT13END

The original scalar measure is

ENDPOINTSOURCE62EECF4C17912CE87EEDAFB1SLOT14END

and the norm is

ENDPOINTSOURCE62EECF4C17912CE87EEDAFB1SLOT15END

For ENDPOINTSOURCE62EECF4C17912CE87EEDAFB1SLOT16END, the canonical matrix and source section remain

ENDPOINTSOURCE62EECF4C17912CE87EEDAFB1SLOT17END

Write ENDPOINTSOURCE62EECF4C17912CE87EEDAFB1SLOT18END. The established cyclic retraction composed with the full jet map gives j\_E on the stated test-function domain, with ENDPOINTSOURCE62EECF4C17912CE87EEDAFB1SLOT19END and ENDPOINTSOURCE62EECF4C17912CE87EEDAFB1SLOT20END. The source relation is

ENDPOINTSOURCE62EECF4C17912CE87EEDAFB1SLOT21END

These are inherited interfaces, not analytic theorems re-proved by the finite tests of this note. The empty packet E=0 has the empty determinant one; all q-positive constructions below concern q\textgreater0.

\subsection{4. SGA 5's relative duality, calculated on this polynomial algebra}\label{endpoint-sga_trace_period--sga-5s-relative-duality-calculated-on-this-polynomial-algebra}

Let R be a commutative ring, let χ be monic of degree q, and put

ENDPOINTSOURCE62EECF4C17912CE87EEDAFB1SLOT22END

The basis 1,S,\ldots,S\^{}(q−1) proves f finite free of degree q. Set

ENDPOINTSOURCE62EECF4C17912CE87EEDAFB1SLOT23END

Here {[}S\^{}(q−1){]} means coefficient extraction in the actual monic remainder. Over ℂ it equals the sum of finite residues of P(S)dS/H(S), with the positive finite-residue orientation. Equivalently it is the coefficient of S\^{}(−1) in P/H, which is the negative of the residue at infinity.

The bilinear form

ENDPOINTSOURCE62EECF4C17912CE87EEDAFB1SLOT24END

is perfect. Its Gram in the unchanged power basis has zeros below total exponent q−1 and ones on that antidiagonal. Thus

ENDPOINTSOURCE62EECF4C17912CE87EEDAFB1SLOT25END

This proves perfection over the original coefficient ring, without dividing by a discriminant. The finite duality functor is

ENDPOINTSOURCE62EECF4C17912CE87EEDAFB1SLOT26END

in degree zero. Equation (14) is its explicit trivialization; evaluation at 1 is its counit. The factor dS/H and the conormal generator {[}H{]} account for the smooth dimension-one and regular codimension-one shifts, which cancel in this finite relative dimension-zero calculation.

\subsubsection{4.1 The actual diagonal class}\label{endpoint-sga_trace_period--the-actual-diagonal-class}

In B\_t=E\_t⊗\_(R{[}t{]})E\_t, with coordinates X,Y, define

ENDPOINTSOURCE62EECF4C17912CE87EEDAFB1SLOT27END

The fraction denotes polynomial division with zero remainder; there is no localization at X−Y. Its defining identities are

ENDPOINTSOURCE62EECF4C17912CE87EEDAFB1SLOT28END

where μ is multiplication E\_t⊗E\_t→E\_t.

The first identity shows the class is supported on the diagonal. Moreover

ENDPOINTSOURCE62EECF4C17912CE87EEDAFB1SLOT29END

To prove it, first take P=1: the coefficient of X\^{}(q−1) in (17) is 1. The identity (X−Y)𝒞\_H=0 then gives (P(X)−P(Y))𝒞\_H=0 for every polynomial P, proving (19). This also proves that 𝒞\_H is the coevaluation tensor inverse to (14).

Taking the categorical trace of multiplication gives the complete composite

ENDPOINTSOURCE62EECF4C17912CE87EEDAFB1SLOT30END

Consequently

ENDPOINTSOURCE62EECF4C17912CE87EEDAFB1SLOT31END

This proves the one-variable instance of the supplied SGA 5 III (6.8.5). In that statement the regular-immersion class is the conormal differential {[}H{]}↦χ'dS. Its trace map is exactly (20), not an additional sign-free identification.

The bilinear trace form is connected to the perfect residue form by the exact map

ENDPOINTSOURCE62EECF4C17912CE87EEDAFB1SLOT32END

This proves the complete radical statement for this finite trace form. At a complex root of multiplicity m its radical has dimension m−1; the perfect residue form itself remains nondegenerate. It is not the assertion of positivity for the original sesquilinear Weil form. Their connection is the retained reflection and Jacobian composite (24).

All maps (12)--(21) commute with coefficient base change: their entries are polynomial expressions in the coefficients and monic remainders. For tensor products, the coevaluation is the external product of the displayed tensors; the diagonal Jacobian is the product of the χ\_i'. If the coefficient object is placed in degree k, its supertrace is (−1)\^{}k times (21). No geometric weight is inferred from that cohomological sign.

\subsubsection{4.2 The source's supported relation maps to this same class}\label{endpoint-sga_trace_period--the-sources-supported-relation-maps-to-this-same-class}

At t=0, the original sum-direction logarithm operator is

ENDPOINTSOURCE62EECF4C17912CE87EEDAFB1SLOT33END

Differentiating the full Mellin amplitude 𝒰(𝐬)P(S), where 𝒰=∏v\_h(s\_i), gives both 𝒰P′ and (∂𝒰)P. Applied to an original relation χP, the latter contains χ and is killed by the original arithmetic quotient. Therefore

ENDPOINTSOURCE62EECF4C17912CE87EEDAFB1SLOT34END

The arithmetic unit is retained in the map before j\_E; it is removed only by the specified inverse/retraction in j\_E. Formula (23) does not differentiate an already-formed scalar e. It differentiates an actual source relation before that relation is sent to its supported quotient zero.

The composition of (23) with λ\_0 is the trace (21). More generally, with the original reflection f\^{}dagger\_k(S)=overline(f(k−overline S)),

ENDPOINTSOURCE62EECF4C17912CE87EEDAFB1SLOT35END

The induced trace in the full arithmetic module still uses η and its retained complement. It is not equated with the cyclic trace without that inclusion and complement.

The conormal lift C{[}S{]}/(χ²) used to compute (23) maps surjectively to the actual contracted level C{[}S{]}/(χ\_2), because (χ²)⊆I²∩C{[}S{]}=(χ\_2). No equality of these ideals is asserted. This keeps the formalization session's correction to higher relation depth.

\subsection{5. The whole derived diagonal is retained}\label{endpoint-sga_trace_period--the-whole-derived-diagonal-is-retained}

There is a complete explicit resolution behind (18). Put d=X−Y and b=𝒞\_H. Over B\_t it is

ENDPOINTSOURCE62EECF4C17912CE87EEDAFB1SLOT36END

\textbf{Exactness.} Work first in T=(R{[}t,Y{]}/(H(Y))){[}X{]}. Then B\_t=T/(db). Both d and b are monic polynomials in X, hence non-zero-divisors in T even when the coefficient algebra has nilpotents. If da∈(db), cancellation of d in T gives a∈(b). If ba∈(db), cancellation of b gives a∈(d). These are precisely the alternating kernel--image identities in (25); B\_t/(d)=E\_t supplies the augmentation.

Pullback to the diagonal is tensoring (25) by μ. The two differentials become 0 and multiplication by χ′. It follows that

ENDPOINTSOURCE62EECF4C17912CE87EEDAFB1SLOT37END

These are module identifications; no multiplication or shuffle structure is silently included. Multiplication by X commutes with every map of (25) and specializes to the original S-action on (26).

At a complex root of multiplicity m, χ′ is the product of m(S−λ)\^{}(m−1) with the actual nonzero local unit. Both positive-degree modules have dimension m−1. For m=1 they vanish. For m\textgreater1 the higher classes are retained as distinct cohomological degrees over the same split scalar support.

This is a classical hypersurface calculation performed here on the exact polynomial. It computes the derived diagonal; it does not assert that these HH degrees equal the separately indexed original nullity ladder. The map linking this calculation to that ladder is the shared conormal differential (23), together with the displayed quotient from the original contracted relation levels.

\subsection{6. The preceding exponential family and its exact coefficient connections}\label{endpoint-sga_trace_period--the-preceding-exponential-family-and-its-exact-coefficient-connections}

Now work over ℂ and write

ENDPOINTSOURCE62EECF4C17912CE87EEDAFB1SLOT38END

The original polynomial is the marked coefficient value 𝐜=𝐜\_h,k. None of its coefficients is shifted or rescaled.

The previous two-term coefficient complex is

ENDPOINTSOURCE62EECF4C17912CE87EEDAFB1SLOT39END

Its unique degree-\textless q remainder gives a free degree-one module. At u=t=0 it is E. The classical finite algebra and this differential quotient have the explicit coefficient-linear inverse maps

ENDPOINTSOURCE62EECF4C17912CE87EEDAFB1SLOT40END

Both normal-form operations fix every degree-\textless q polynomial, proving the inverse laws. The map is a coefficient-module isomorphism, not a declaration that multiplication descends through D. For u≠0 its connections are

ENDPOINTSOURCE62EECF4C17912CE87EEDAFB1SLOT41END

The commutators with D vanish term by term. The quotient is by the linear range of D, not an algebra ideal: {[}S(χ−t){]}=−u{[}1{]}. The previous note's source deformation

ENDPOINTSOURCE62EECF4C17912CE87EEDAFB1SLOT42END

satisfies j\_ED\_t\^{}(k)=A(t)j\_E and D\_t\^{}(k)r\_N−r\_NA(t)=dK\_N\^{}prim, with A(t)=A+t𝓡. Its original source relation is unchanged.

Let C\_a be the degree-\textless q matrix for reduction of S\^{}(a+1) times a representative by D/dS, and let A(t) be the usual companion matrix for χ−t. The period matrix obeys

ENDPOINTSOURCE62EECF4C17912CE87EEDAFB1SLOT43END

\subsubsection{6.1 The trace of the reduced multiplication is exactly classical}\label{endpoint-sga_trace_period--the-trace-of-the-reduced-multiplication-is-exactly-classical}

Although C\_a need not equal A(t)\^{}(a+1),

ENDPOINTSOURCE62EECF4C17912CE87EEDAFB1SLOT44END

Here is the explicit comparison. Reducing S\^{}(b+a+1), with 0≤b≤q−1 and 0≤a≤q−1, starts with degree at most 2q−1. Every term introduced by the uP′ part has degree at most b+a−q≤b−1. Later reductions only lower degree. Thus C\_a−A(t)\^{}(a+1) is strictly upper triangular in the fixed increasing-power frame: its row indices are strictly below its column index. Its diagonal entries vanish, proving (32). No u-dependent coefficient is discarded from the full matrix.

\subsection{7. Complete evaluation of the determinant of periods}\label{endpoint-sga_trace_period--complete-evaluation-of-the-determinant-of-periods}

Set d=q+1. Fix u≠0 and a choice of arg u on a sector. Retain the rays

ENDPOINTSOURCE62EECF4C17912CE87EEDAFB1SLOT45END

The periods are

ENDPOINTSOURCE62EECF4C17912CE87EEDAFB1SLOT46END

The leading exponent on each ray is −r\^{}d/(d\textbar u\textbar). Lower-degree terms are uniformly dominated on compact coefficient sets. This proves convergence and all coefficient differentiations in (31). Integrating D(P)e\^{}(Φ\_t/u) gives a total derivative whose endpoints at infinity vanish and whose two origin values cancel with the stated orientations.

Define the polynomial

ENDPOINTSOURCE62EECF4C17912CE87EEDAFB1SLOT47END

Both descriptions are finite coefficient calculations. At a complex fibre it is the sum of all critical values Φ\_t(λ), counted with the multiplicity of each root of χ−t.

Its gradients are

ENDPOINTSOURCE62EECF4C17912CE87EEDAFB1SLOT48END

\textbf{Proof.} On the open set of simple roots α\_j, differentiate ∑Φ\_t(α\_j). The α\_j derivatives are multiplied by Φ\_t′(α\_j)=χ(α\_j)−t=0. The remaining terms are exactly (36). Both sides of (36) are polynomials in the coefficients, obtained by companion traces. Their equality on the nonempty simple-root open set proves the polynomial identity everywhere, including all collisions. This extension does not replace the fibres by their reductions: their ranks and traces retain multiplicities.

Equations (31), (32), and Jacobi's determinant derivative give

ENDPOINTSOURCE62EECF4C17912CE87EEDAFB1SLOT49END

These determinant equations hold even before invertibility is known, by the adjugate identity. At the monomial member 𝐜=0,t=0, 𝔽=0 and direct integration gives

ENDPOINTSOURCE62EECF4C17912CE87EEDAFB1SLOT50END

The last matrix F\_d has Gram F\_d\^{}*F\_d=d(I\_q+J\_q), where J\_q is the all-ones matrix. Indeed expansion of the finite geometric sums gives diagonal 2d and off-diagonal d.~Hence \textbar detF\_d\textbar²=d\^{}(q+1), in particular it is nonzero.

Along the straight coefficient path from the monomial member to any chosen (𝐜,t), (37) is a scalar differential equation. Integrating it with its actual initial value proves

ENDPOINTSOURCE62EECF4C17912CE87EEDAFB1SLOT51END

There is no path-dependent logarithm in this equation. In particular the period matrix is invertible for every u≠0, at repeated as well as simple critical points.

The complete unsquared constant is retained by (38), its determinant, the ray ordering, and the choice of arg u. Taking squared moduli and using the gamma multiplication identity

ENDPOINTSOURCE62EECF4C17912CE87EEDAFB1SLOT52END

gives

ENDPOINTSOURCE62EECF4C17912CE87EEDAFB1SLOT53END

The factor (2π\textbar u\textbar)\^{}q is calculated from (38)--(40); it is not assigned a value by choosing a new basis. Formula (38) remains the authoritative phase-containing version.

This explicitly finishes the determinant calculation for the previously constructed exponential comparison. It is consistent with the classical exp(f) determinant theory; it is not a claim that the full period matrix, or its arithmetic condition number, is determined by its determinant.

\subsection{8. Recover the complete cyclic spectrum from that determinant family}\label{endpoint-sga_trace_period--recover-the-complete-cyclic-spectrum-from-that-determinant-family}

For 1≤r≤q, equation (37) gives the logarithmic derivative without choosing a logarithm branch:

ENDPOINTSOURCE62EECF4C17912CE87EEDAFB1SLOT54END

Put d\_0=1 and recursively define

ENDPOINTSOURCE62EECF4C17912CE87EEDAFB1SLOT55END

Then

ENDPOINTSOURCE62EECF4C17912CE87EEDAFB1SLOT56END

\textbf{Proof.} In the formal variable z,

ENDPOINTSOURCE62EECF4C17912CE87EEDAFB1SLOT57END

Coefficient extraction gives (43), and the determinant is a degree-q polynomial. Because A(t) is the specified cyclic companion action, its characteristic polynomial is also its minimal polynomial. Thus this recovery retains one complete length-m block at every root of multiplicity m. For a general noncyclic larger arithmetic module, traces do not classify all Jordan blocks; its inclusion η and complementary module stay in the source programme.

The precise map chain is

ENDPOINTSOURCE62EECF4C17912CE87EEDAFB1SLOT58END

Only the full coefficient-family derivatives in (42), not a single determinant value, give the recovery.

\subsection{9. Two explicit calibrations}\label{endpoint-sga_trace_period--two-explicit-calibrations}

For χ(S)=S−λ,

ENDPOINTSOURCE62EECF4C17912CE87EEDAFB1SLOT59END

For u\textgreater0, with the original contour oriented upwards on iℝ,

ENDPOINTSOURCE62EECF4C17912CE87EEDAFB1SLOT60END

In the ordering (33), q=1 gives Γ₁ oriented downwards, so ENDPOINTSOURCE62EECF4C17912CE87EEDAFB1SLOT61END and the period in (48) is the negative of the sole entry of (34). This is the exact orientation map; the squared determinant is unchanged. The full nonzero λ remains; a cohomological dimension or pure finite-field weight does not set it to a critical-line value.

For χ(S)=S²+c\_1S+c\_0,

ENDPOINTSOURCE62EECF4C17912CE87EEDAFB1SLOT62END

At the discriminant c\_1²−4(c\_0−t)=0 the formula stays unchanged and the fibre retains its double-root algebra. This is tested symbolically. Independent ray quadrature for small declared complex coefficient fixtures is reported separately, without an interval-certification claim.

\subsection{10. Transport the original metric and control through the same periods}\label{endpoint-sga_trace_period--transport-the-original-metric-and-control-through-the-same-periods}

At the actual marked coefficients, use the original arithmetic G\_N of (10). On the period-coordinate target define the \textbf{transported original metric}

ENDPOINTSOURCE62EECF4C17912CE87EEDAFB1SLOT63END

The map Π:(E,G\_N)→(ℂ\^{}q,ENDPOINTSOURCE62EECF4C17912CE87EEDAFB1SLOT64END) is an isometry by this explicit formula, not by an assignment of the original metric to I\_q. Conversely Π\^{}\emph{Π is the Gram for the separately specified coordinate metric I\_q, and its comparison operator is G\_N\^{}(−1)Π\^{}}Π.

The exact determinant formula is

ENDPOINTSOURCE62EECF4C17912CE87EEDAFB1SLOT65END

At a common (𝐜,u,t), the complete marked endpoint ratio satisfies

ENDPOINTSOURCE62EECF4C17912CE87EEDAFB1SLOT66END

Thus every common period determinant factor in that ratio is now explicitly known and cancels by (51), with no remaining unknown proportionality constant. This identifies exactly which information the period determinant does and does not add to the endpoint estimate.

For fixed u\textgreater0 and real t, differentiate (50), retaining the original G\_N fixed and using Π′=−ΠA(t)/u. One obtains

ENDPOINTSOURCE62EECF4C17912CE87EEDAFB1SLOT67END

and therefore

ENDPOINTSOURCE62EECF4C17912CE87EEDAFB1SLOT68END

At ENDPOINTSOURCE62EECF4C17912CE87EEDAFB1SLOT69END, ENDPOINTSOURCE62EECF4C17912CE87EEDAFB1SLOT70END is the original arithmetic control. For any vector x, the coordinate change ENDPOINTSOURCE62EECF4C17912CE87EEDAFB1SLOT71END carries both numerator and denominator of its generalized Rayleigh quotient exactly. Equations (30) and (11) show that the source theta primitive remains the same along this specific t-deformation.

This is a differential-metric reformulation of the same problem. The finite-field pure cohomology of the auxiliary family controls its Frobenius, not automatically the right side of (54). The required relation between those operators has not been supplied by taking a determinant.

\subsubsection{10.1 The constituent period minor detects the arithmetic defect}\label{endpoint-sga_trace_period--the-constituent-period-minor-detects-the-arithmetic-defect}

The full determinant is sensitive to all coefficient derivatives by (42), but along the original t-direction it sees only Tr A. In a reflection-stable arithmetic packet of dimension q, ENDPOINTSOURCE62EECF4C17912CE87EEDAFB1SLOT72END. For the full conjugation-stable quartet packet, ENDPOINTSOURCE62EECF4C17912CE87EEDAFB1SLOT73END exactly. This is already the total weight balance; an off-line pair can retain that same real trace.

The correct retained subobject is an actual ENDPOINTSOURCE62EECF4C17912CE87EEDAFB1SLOT74END-invariant inclusion ENDPOINTSOURCE62EECF4C17912CE87EEDAFB1SLOT75END, ENDPOINTSOURCE62EECF4C17912CE87EEDAFB1SLOT76END. Let ENDPOINTSOURCE62EECF4C17912CE87EEDAFB1SLOT77END and define the \textbf{rectangular} period map and its two Grams

ENDPOINTSOURCE62EECF4C17912CE87EEDAFB1SLOT78END

The determinant here is the norm of the actual exterior map ENDPOINTSOURCE62EECF4C17912CE87EEDAFB1SLOT79END on ENDPOINTSOURCE62EECF4C17912CE87EEDAFB1SLOT80END; no determinant of a rectangular matrix is taken. The domain Gram retains the theta-source inclusion ENDPOINTSOURCE62EECF4C17912CE87EEDAFB1SLOT81END and all its jets. For real ENDPOINTSOURCE62EECF4C17912CE87EEDAFB1SLOT82END and ENDPOINTSOURCE62EECF4C17912CE87EEDAFB1SLOT83END,

ENDPOINTSOURCE62EECF4C17912CE87EEDAFB1SLOT84END

This follows by differentiating the displayed Gram and taking its trace. It is not a purity assumption on F. For the full positive-real-defect generalized eigenspace ENDPOINTSOURCE62EECF4C17912CE87EEDAFB1SLOT85END, the right side is the exact aggregate ENDPOINTSOURCE62EECF4C17912CE87EEDAFB1SLOT86END already used by the exterior theorem. Thus the period family has an explicit observable for the original off-line defect, not just for its already-balanced total trace.

The original extension to the other constituents is present as well. Let P be the orthogonal projection in the specified period-coordinate metric onto ENDPOINTSOURCE62EECF4C17912CE87EEDAFB1SLOT87END. Since ENDPOINTSOURCE62EECF4C17912CE87EEDAFB1SLOT88END,

ENDPOINTSOURCE62EECF4C17912CE87EEDAFB1SLOT89END

Both identities are direct differentiations. The first derivative is internal to F; the first possible normal acceleration is the exact rank-at-most-one map on the right. Its kernel, image, and norm remain available. A repeated-root primary subspace is allowed; no eigenspace basis or semisimple replacement is used.

Equations (54a)--(54c) specify the constituent period object that can be estimated next. Formula (41) finishes the full determinant, but does not by itself bound these minor norms or their comparison with ENDPOINTSOURCE62EECF4C17912CE87EEDAFB1SLOT90END. The preceding logarithmic-constituent extension and its coupling subtraction remain the source-specific metric input.

\subsection{11. The source-specific remaining term stays explicit}\label{endpoint-sga_trace_period--the-source-specific-remaining-term-stays-explicit}

For the current programme, the original endpoint cost is

ENDPOINTSOURCE62EECF4C17912CE87EEDAFB1SLOT91END

The preceding logarithmic-constituent calculation already supplies a genuine source upper certificate. For an A-invariant inclusion I:F→E and its G\_j-orthogonal coefficient lift H\_0 of E/F, let C=G\_i−G\_j and

ENDPOINTSOURCE62EECF4C17912CE87EEDAFB1SLOT92END

Put Q\_+=Q\_0+C\_Q and

ENDPOINTSOURCE62EECF4C17912CE87EEDAFB1SLOT93END

That construction proves the exact decomposition

ENDPOINTSOURCE62EECF4C17912CE87EEDAFB1SLOT94END

These are the source matrices that still require uniform control. Under the coordinate change Π, all inclusions, quotient lifts, metrics and D\_F are transported together; (58) retains the same values. A period basis does not eliminate the coupling or the theta mass.

The current GitHub statement retains the already derived off-line-quartet consequence

ENDPOINTSOURCE62EECF4C17912CE87EEDAFB1SLOT95END

This note does not establish a contradictory upper estimate. It finishes the diagonal trace and the determinant-period calculation, recovers all cyclic power traces from that determinant family, and gives the exact unchanged theta-metric comparison (51)--(54). The quantitative focus remains a bound on (58), or equivalently on the same original generalized control in (54), rather than another uncomputed determinant or an assumption that local support entails metric purity.

\subsection{12. Relation to Deligne's finite-field family}\label{endpoint-sga_trace_period--relation-to-delignes-finite-field-family}

The previous construction takes the finitely generated coefficient ring R\_0 containing the actual c\_a and 1/(q+1)!, and for a specified finite-field specialization of characteristic p\textgreater q+1 obtains

ENDPOINTSOURCE62EECF4C17912CE87EEDAFB1SLOT96END

Its infinity chart is the original one: with w=1/S,

ENDPOINTSOURCE62EECF4C17912CE87EEDAFB1SLOT97END

The derivative in v has the specified inverse v/d.~This is the retained local model for applying Deligne's exponential-sum/lissity theorem over u≠0. The generic polynomial has rank-q degree-one compact cohomology and weight-one Frobenius under those hypotheses. Those are the inherited cited Deligne results, not fresh analytic conclusions of this note.

The maps are the span

ENDPOINTSOURCE62EECF4C17912CE87EEDAFB1SLOT98END

The SGA trace formalism applies in each specified coefficient realization with its own endomorphism. It does not supply a homomorphism ℂ→𝔽\_Q or an intertwiner from the auxiliary Frobenius to the original A. The actual complex source/operator comparison is (30), (31), and (54). This keeps the two realizations connected at their common coefficient model without filling the missing operator comparison with an unsupported equality.

\subsection{13. Verification and reproducibility}\label{endpoint-sga_trace_period--verification-and-reproducibility}

\texttt{check\_trace\_period.py} contains eighteen finite exact regression methods: residue duality, the diagonal tensor, the trace/Jacobian composite, repeated-root fibres, periodic diagonal resolution, coefficient-base-change identities, quantum-reduction traces, critical-value gradients, Newton recovery, source conormal response, full metric congruences, and split quotient labels. Separate negative controls alter the Jacobian, discard the quantum derivative, or erase multiplicity; each must fail in both normal and optimized Python.

The period convergence, determinant formula for arbitrary q, and all-general resolution exactness have written proofs above. Finite symbolic or numerical tests are not replacements for those proofs. \texttt{period\_numerics.py} performs additional explicitly non-certified contour comparisons on declared small fixtures. No new Lean execution, arithmetic zero location, interval certificate, or uniform purity theorem is claimed.

The package does not redistribute the uploaded SGA corpus. The source inventory retains hashes, original paths and read ranges. The GitHub read was pinned to c05c709d707bacda335deff70fbb5b841a751d1f. No remote file is changed by building this package.

\subsection{References and source status}\label{endpoint-sga_trace_period--references-and-source-status}

\begin{enumerate}
\def\labelenumi{\arabic{enumi}.}
\tightlist
\item
  A. Grothendieck, written up by L. Illusie, \emph{SGA 5, Exposé III: Lefschetz Formula}, §§4.1, 4.7--4.8, and coherent appendix §§6.1--6.8, especially (6.8.5). Basis: the user's supplied English master, read directly as LaTeX. Source wording and qualifications are recorded separately.
\item
  P. Deligne, \emph{La conjecture de Weil II}, Publications Mathématiques de l'IHÉS 52 (1980), 137--252. The previous source-reviewed construction uses §§3.7.2--3.7.4. This execution reads that predecessor's TeX and preserves its stated scope, rather than claiming a new complete scan verification.
\item
  S. Bloch and H. Esnault, \emph{Gauß--Manin determinant connections and periods for irregular connections}, arXiv:math/9912095 (1999). The abstract was inspected for attribution; its TeX download failed, so no detailed theorem of that article is asserted as independently read here.
\item
  NIST Digital Library of Mathematical Functions, §5.5(iii), gamma multiplication formula. It is used only to evaluate the explicit product in (38); the full product and phases remain available.
\item
  Parent programme sources: \emph{Tau\_Deligne\_Exponential\_Comparison} and \emph{Tau\_Deligne\_Logarithmic\_Constituent\_Control}, 13 September 2026. Their original sections, arithmetic source maps, metric definitions, and verification scopes are retained. General source and analytic constructions are inherited where expressly marked above.
\end{enumerate}
